\section{総括 — 5月、生成AIの競争軸はモデルから実行環境へ広がった}
\label{sec:retrospective-synthesis}
\noindent\textbf{本号の各章を重ねると、2026年5月は単なるmodel releaseの月ではなかった。default modelとAPI lifecycleが動く一方で、agentを置くsandboxとnetwork policy、長時間workloadを支えるserving stack、memory／retrievalを守るsecurity、そしてbenchmark外のcapabilityを測るevaluation設計まで、生成AIの実行条件そのものが競争軸として前景化した。}
\medskip
\begin{multicols}{2}
\subsection{Frontier Models — model名ではなくdelivery surfaceを見る}
GPT-5.5 Instant、AnthropicのMay release、Gemini APIを並べると、5月の更新はmodel identityだけでなくdefault experience、GA／preview、agent integrationといったdelivery surfaceに現れた。利用者にとっては性能表だけでなく、いつ何がどの提供形態で変わるかがarchitecture上の条件になる。 \cite{src-01a8b52eee81,src-f6aa679babd7,src-a24fd97eb3d0}

\subsection{Agents \& Runtime — capabilityを実行してよい範囲を設計する}
Running Codex safelyとWindows sandbox、Kimi Code CLI、MistralのMCP／workflow更新を接続すると、agent engineeringの中心はtoolを増やすことだけではない。file、network、approval、telemetry、CLI／protocolを通じて、長時間動作をどの境界の中へ置くかが実運用の設計対象になった。 \cite{src-66be12207714,src-4a3442a5ca8c,src-9191d43918fc,src-d8b7c1c36830}

\subsection{Inference \& Serving — agent時代のworkloadをstack全体で支える}
vLLM、SGLang、FlashInferのMay releasesでは、attention backend、cache、quantization、MoE、model-specific pathが継続的に更新された。agent workloadが複数callと成長するcontextを持つほど、model qualityだけでなくruntimeが長時間taskのcostとlatencyを決める。 \cite{src-77e29e3fac8f,src-c97320147e8c,src-80b8d67933c0,src-5a693cfc6852,src-69c5523286bd,src-74a7052acd7a,src-ac04c00cf21f}

\subsection{Safety \& Security — 長時間stateは新しいtrust boundaryになる}
ClawTrojanとMemPoisonは異なるthreat modelだが、どちらも一回のpromptを超えてstateへ残るcontrolを問題にする。workspace、long-term memory、retrievalのprovenanceを追跡しなければ、agentの安全性をper-request filterだけで閉じることは難しくなる。 \cite{src-7fa0ea6bcf47,src-f216d6cd7707,src-062c2718dd96}

\subsection{Capability \& Evaluation — 成果と評価方法を同時に読む}
離散幾何の事例はbenchmark外のresearch taskでmodel capabilityが現れる可能性を示すfirst-party reportである一方、third-party evaluation playbookはcapability、safeguard、evaluation validityをどう切り分けるかを問う。単一の成功例を一般能力へ広げず、評価手続きそのものを証拠の一部として扱う必要がある。 \cite{src-1c93d6c1c725,src-8c84b8422093,src-24add4f36a45}

\end{multicols}
\Needspace{0.34\textheight}
\sectionkicker{Cross-chapter synthesis}
\subsection*{各章はどうつながるか}
\addcontentsline{toc}{subsection}{各章はどうつながるか}
\begin{center}
\small
\renewcommand{\arraystretch}{1.16}
\begin{tabularx}{\linewidth}{@{}p{0.26\linewidth}X@{}}
\toprule
接続する章 & 本号で見える構造的な関係 \\
\midrule
Frontier Models $\rightarrow$ Agents \& Runtime & model／APIの更新がagent surfaceへ入るほど、capabilityと実行権限を別々に設計する必要が強まる。 \cite{src-a24fd97eb3d0,src-66be12207714} \\
Agents \& Runtime $\rightarrow$ Safety \& Security & 長時間動作とpersistent stateは、workspaceやmemoryを新しいtrust boundaryへ変える。 \cite{src-4a3442a5ca8c,src-f216d6cd7707} \\
Inference \& Serving $\rightarrow$ Agents \& Runtime & agent workloadの複数call・成長contextは、cacheやruntime最適化をmodel外の主要設計問題にする。 \cite{src-ac04c00cf21f,src-77e29e3fac8f,src-c97320147e8c,src-80b8d67933c0} \\
Capability $\rightarrow$ Evaluation & benchmark外の成果が増えるほど、第三者評価とmeasurement validityの設計がcapability claimの解釈に不可欠になる。 \cite{src-1c93d6c1c725,src-8c84b8422093,src-24add4f36a45} \\
\bottomrule
\end{tabularx}
\renewcommand{\arraystretch}{1.0}
\end{center}
\normalsize
\begin{claimboundary}[この総括の境界]
この総括は本号で選択した一次資料と原論文を横断した編集上のsynthesisである。vendor／project／authorが報告した性能、attack success、capability比較を独立再現済みの事実へ昇格させず、異なる評価条件の数値を直接比較しない。
\end{claimboundary}
\subsection*{5月をどう位置づけるか}
\addcontentsline{toc}{subsection}{5月をどう位置づけるか}
2026年5月は、より賢いmodelが登場した月であると同時に、その能力を現実のworkflowへ持ち込むためのexecution environmentが急速に具体化した月だった。6月以降にagentのmemory、scheduling、authorization、scientific workflowがさらに前景化する前段として、5月はsandbox・serving・persistent state・evaluationという基盤が同時に動き始めた転換点として位置付けられる。 \cite{src-01a8b52eee81,src-66be12207714,src-77e29e3fac8f,src-c97320147e8c,src-80b8d67933c0,src-f216d6cd7707,src-8c84b8422093}
